\section{Logical rank and the no-conversion checking proof}
\label{sec:annex-p07-ch88-checking}

\begin{annexitem}{annex:p7.ch88.check-generation}{thm:p07-ch88-check-generation}{Actual source environments and last-rule generation}
\begin{proof}
Induct jointly on validity, typing and equality derivations, retaining a
certificate for every older prefix of the environment. Empty validity has
the empty prefix. Logical extension retains the prefix certificates from
its validity premise and appends that very annotation certificate; object
extension does so with its type-\(\ast\) premise. A logical-product,
object-product or abstraction premise is in an extended environment:
its induction already provides all its prefix certificates, so remove
the last declaration to obtain the conclusion's prefix data. Projection
has its explicit validity premise. Application and typing conversion
retain the common environment data from their typing premise. Equality
reflexivity has a validity or typing premise. Symmetry and transitivity
retain their premise data. Product equality has the annotation equality
premise in the conclusion's context; lambda/application equality also
have their stated typing premises. Beta equality has the body premise
in an extended valid environment, whose prefix data again suffice.
The restricted guarded cases are the same cases with additional
premises, not an extra environment rule. This proves the asserted
regularity without typing all raw equality endpoints. Each annotation
certificate can be retained as an actual subderivation: it is supplied
at the extension introducing that declaration, and otherwise inherited
from premise prefix data. A binder's annotation certificate has smaller
height than the judgment using its binder premise.

For generation induct on the typing derivation. A final typing
conversion leaves subject and environment literally unchanged: append
its displayed source equality to the chain already obtained from the
premise. Otherwise the last common rule is forced by the subject's
literal constructor. Its actual premises and conclusion give precisely
the asserted initial type. No other common rule has that subject shape,
and there is no independent source weakening rule to hide another case.
The chain is empty when no conversion was stripped. In particular this
does not invent reflexivity at a potentially unsorted initial type.
\end{proof}
\end{annexitem}

\begin{annexitem}{annex:p7.ch88.check-substitution}{thm:p07-ch88-check-substitution}{Thinning and substitution before subject reduction}
\begin{proof}
Use the proved named binding correspondence for these rule inductions,
then decode to the actual indexed vectors. Names inserted into a
context or bound locally are fresh. This uses only the established
raw relocation, lookup and substitution identities, not a typing inverse.

For thinning induct on derivation height, simultaneously for all
three judgments and all valid larger environments retaining the
old declarations in order. A star-validity conclusion is the larger
environment's assumed validity. Projection has the same declaration
there, with its relocated lookup. In a binder rule the preceding
regularity result supplies its annotation sort certificate as an
actual smaller subderivation. Thin that certificate first, extend
the larger context by the same annotation, and apply induction to
the body premise in that valid extension. Rebuild the logical
product, object product or abstraction. Application thins both
premises; typing conversion thins the typing and equality premises.
Equality reflexivity, symmetry and transitivity use their premise
inductions. Product equality uses the first annotation equality
and its body equality in the valid first-domain extension. Lambda
equality retains its body typing there. Application equality retains
its stated application typing and one or both child equalities.
Beta equality retains its body typing in the extension and actual
argument typing in the larger base. In the restricted cases the
logical guards remain logical because relocation preserves source
classes. These are all primitive rules. Fresh binder renaming avoids
capture, and indexed decoding gives the stated relocated judgment.

For substitution induct on the given derivation's height, simultaneously
for all judgments and suffixes. First construct the substituted suffix
oldest-first. Its annotation certificates are smaller actual
subderivations by regularity. Apply substitution induction to the
first with empty suffix, and then to each next certificate with
the already constructed preceding substituted suffix. The resulting
object-type or logical-context certificate licenses each new
declaration. Thus the target context is valid before the rule case
requires projection or a binder extension.

At projection of the removed variable, insert the given \(N:A\)
thinned to this valid target suffix. Its type \(A\), scoped over
the older \(\Gamma\), contains no removed reference. At any other
projection, its substituted annotation is exactly its declaration
in the new suffix. Validity star uses that new context validity;
extension uses its annotation induction. Logical/object products
and abstraction use body induction in their substituted extensions.
Application combines its two inductions and the substitution
composition identity. Typing conversion combines its two premises.
Equality reflexivity, symmetry and transitivity combine their
corresponding inductions; product/lambda congruence retains the
first-annotation body context and the actual typing guard. Application
congruence retains its application typing and child equalities.
For beta equality the substituted premises are precisely primitive
beta premises, with the new contractum identified by composition.

That identity is a structural induction in the two raw source classes.
At a variable distinguish the selected name, a fresh locally bound
name, and every other name; the inserted argument is scoped over
\(\Gamma\), so has no removed reference. Application distributes
substitution. For a binder choose its name fresh for both arguments,
apply the annotation induction at the old scope and the body induction
at the extended scope. Star is fixed. The resulting named identity
is literal after indexed decoding. Object substitution preserves
the object/logical classes, so every restricted logical guard
survives. No case invoked type uniqueness, beta subject reduction,
normalization or equality reflection from a target calculus.
\end{proof}
\end{annexitem}

\begin{annexitem}{annex:p7.ch88.check-rank}{thm:p07-ch88-check-rank}{Rank of source classes and substitution}
\begin{proof}
Every object has rank zero, regardless of its internal object syntax.
Substitution of an object for an index stays in the object class by
\cref{thm:p07-ch88-index-operations}. For a logical context, induct on
its grammar. The star stays star. A product has annotation in one of
the two source classes and logical-context body. An object annotation
stays an object of rank zero; a logical annotation uses the simultaneous
logical induction; the body uses that induction at the increased
substitution cutoff. The maximum defining the product rank is unchanged.
Relocation is covered by the same argument. This establishes rank
stability under the actual object-argument substitution.

Induct on source equality. The two reflexivities preserve class and
rank; symmetry reverses that assertion; transitivity composes it.
Lambda and application equalities have object endpoints, hence rank
zero. Beta has an object redex and an object contractum by scoped
substitution, again both rank zero. In product congruence the body
induction gives the same body class on both sides. If it is object,
both entire products are objects of rank zero. If it is logical,
both products are logical and the two annotation ranks and body ranks
agree by their inductions; the maxima consequently agree. All full
and restricted guarded replacements fit these cases. This proves the
printed rank-invariance conclusion, but says nothing about newly
created redexes in an enclosing application.
\end{proof}
\end{annexitem}

\begin{annexitem}{annex:p7.ch88.check-domain-rank}{thm:p07-ch88-check-domain-rank}{The increasing innermost example with all guards}
\begin{proof}
The named environment is valid by extension first with the logical
annotation \(\ast\), then with \(A:\ast\), then with the logical
annotation \(\ast\). In its fixed indexed form it is
\([A:\ast][a:1][C:\ast]\ast\), so its references to \(A,a,C\)
are respectively \(3,2,1\). Projection gives \(A:\ast\),
\(a:A\), and \(C:\ast\). In the further extension by \(x:A\)
and then \(B:\ast\), projection gives \(B:\ast\). Abstraction
therefore gives \((\lambda B:\ast)B:[B:\ast]\ast\) in the
\(x\)-extended context, and the outer abstraction gives
\(L:[x:A][B:\ast]\ast\). Application to \(a:A\) yields
\(La:[B:\ast]\ast\); the body does not contain \(x\), so the
binder-removal substitution has exactly this result. Application to
\(C:\ast\) yields \(M:\ast\). Independently the single lambda
\((\lambda B:\ast)B\) has type \([B:\ast]\ast\) and hence
\(M':\ast\). These use only the common rules, so both variants
contain the derivations.

The logical beta generator applies to \(La\): its body is the
constructor \((\lambda B:\ast)B\) typed at the logical context
\([B:\ast]\ast\), and its argument has the literal domain type
\(A\). Its contractum is that inner lambda. Logical function
compatibility lifts this edge through application to \(C\), since
the whole application has logical result type \(\ast\). This is
the claimed directed step.

The only beta redex initially is \(La\). Its body is a lambda with
variable body; its annotations are \(A,\ast\); its argument is the
variable \(a\). Thus no inner logical or ordinary beta redex exists,
so it is innermost even under the stronger interpretation which
examines every subfield. Its printed domain complexity is
\(\delta(A)=0\). After the contraction the only redex is the whole
\(M'\), with domain \(\ast\) and complexity one. There were no
other complexities to compensate for this increase. A multiset
extension of the usual strict order on natural ranks cannot decrease
from singleton zero to singleton one. This disproves that precise
measure claim, not normalization or decidability itself.
\end{proof}
\end{annexitem}

\begin{annexitem}{annex:p7.ch88.check-raw-reflection}{thm:p07-ch88-check-raw-reflection}{Reflection only on the actual raw image}
\begin{proof}
Induct on the position of a target step starting at \(TE\). At a
root redex, injectivity of the source constructors makes its function
a source lambda with object body, and its argument a source object.
The already proved substitution identity makes the target contractum
exactly \(T([N/x]M)\); scoped source substitution preserves its
class. At an application child use the induction and rebuild the same
source application. At a product or lambda field use the induction at
the actual old depth for the annotation and increased depth for the
body. These exhaust the target positions; the source image contains
no target-only constant \(\Box\). Thus every step has an actual
source preimage and remains in that image. Repetition gives path
reflection, without decoding arbitrary off-image target syntax.

For two source paths from the same scoped input, their target paths
have a common reduct by the supplied raw annotated CC confluence
(\cref{thm:p12-har-pseudo-confluence}, instantiated to its two sorts).
Both paths from the common initial image to this reduct reflect as
just proved. Injectivity gives a common indexed source reduct. This
establishes source confluence on its raw scoped grammar. For an input
with a CC typing certificate, exact CC strong normalization
\cref{thm:p12-har-cc-normalization} rules out an infinite reflected
path. A fixed strategy selecting the first redex in a finite
depth-first constructor walk therefore terminates at a normal form.
Confluence makes any two such normal forms identical after indexed
decoding. All annotation and type fields are included; no erased
program result is used.
\end{proof}
\end{annexitem}

\begin{annexitem}{annex:p7.ch88.check-normal-transport}{thm:p07-ch88-check-normal-transport}{Direct transport of every p112 rule}
\begin{proof}
Induct on the normal-variant derivation. Use only exact CC's supplied
structural prerequisites: valid contexts/projection/thinning,
generation, type correctness, substitution, context conversion,
uniqueness and full compatible subject reduction, proved in paragraph
(a) of \cref{annex:p12.har.cc-normalization}; raw confluence and
product injectivity accompany them. Its axiom is \(\ast:\Box\)
and its product triples are precisely the four stated there.

Empty validity maps to that axiom. For normalized logical extension,
the induction gives \(T\Delta:\Box\). Strong normalization gives
\(\mathcal N(\Delta)\); reflection identifies its target with the
target normal form. Subject reduction retains type \(\Box\), so
this is a valid target annotation certificate. Extend the context
and thin the axiom there. For normalized object-type extension use
the same reasoning at type \(\ast\). Both conclusions' declaration
annotations are normal by construction.

For a logical product, target validity of the extended premise exposes
an annotation sort \(TA:s_1\). Its body induction has type \(\Box\).
The admitted triple \((s_1,\Box,\Box)\) forms its image product.
For an object product the body induction has type \(\ast\) and
\((s_1,\ast,\ast)\) forms its image at \(\ast\). Projection
uses the actual encoded lookup in the valid target context.

For abstraction, its body induction is \(TM:TP\). The type has a
sort by target type correctness: \(TP\) is a source image and so
is not the exceptional literal \(\Box\); a literal \(\ast\)
has its axiom sort. Combine that sort with the annotation certificate
from extended-context validity. The corresponding one of the four
triples forms the product, and target abstraction applies.

For application, its two inductions yield the literal function product
and argument judgments. Target application types the source
binder-removal instance \(T([N/x]B)\). Target type correctness sorts
this result type; again a source image is not the exceptional top.
Therefore strong normalization applies to the result type itself.
Normalize it, retaining its sort by subject reduction, and convert
the application's type to that normal result. Reflection supplies
the precise source normal form on the sheet, not an erased or
weak-head replacement. This proves availability of the normalizer
at this rule before using its conclusion.

Source context regularity is obtained by the same prefix-carrying
induction as for the common rules, with normalized extensions. Its
validity derivation forces each actual annotation normal. In products
and lambdas the premise environment includes their literal domain,
so that domain is normal. Recursive premises give annotation-normal
bodies and application children. Lookup relocation creates no beta
redex and preserves a normal annotation. The product conclusion has
normal domain and normal body type, hence is normal; the object
product's type is \(\ast\); application normalizes its result.
Induction therefore proves the claimed subject, context and result
normality invariants. No inverse to source typing or full-source
subject reduction entered this proof.
\end{proof}
\end{annexitem}

\begin{annexitem}{annex:p7.ch88.check-normal-complete}{thm:p07-ch88-check-normal-complete}{Successful certificates, structural completeness and finite failure}
\begin{proof}
First work in an already checked context and prove soundness of the
three structural operations by induction on their original expression
inputs. At \(\mathrm{Sort}\), the finite normality test is explicit;
the subsequent validity or type-\(\ast\) derivation gives the
appropriate normal-extension rule, since \(\mathcal N(A)=A\).
\(\mathrm{Valid}\) uses empty validity or the logical-product rule
on its recursive body. Variable inference uses its actual lookup.
An object-product success has the required body type \(\ast\).
A lambda success has the actual body type and its product conclusion.
An application success has literal product function type and literally
matching argument type, hence precisely the two p.~112 premises.
Its result normalizer is defined: the recursively constructed premise
derivations have the CC certificates of the preceding transport proof,
whose application case sorts this result type before normalization.
The resulting source normal form is exactly the rule's conclusion.
All returned types are normal by the transport invariant. Thus every
success builds a genuine normal-variant derivation and a CC certificate.

Prove the stronger completeness assertions simultaneously for an
annotation-normal object with any given CC type, and a logical context
with any given CC type. The latter operation returns validity, and
its target type is beta-equal to \(\Box\). Induct on the source
input syntax, retaining the checked environment, and use the supplied
CC generation theorem for the given typing. For a logical star,
generation exposes its axiom and its validity is immediate. For a
logical product, generation exposes domain/body sort premises. Apply
the smaller annotation induction when its annotation is an object,
or the smaller logical induction when it is a logical context.
An object annotation's returned normal type is beta-equal to its
generated sort. It cannot equal \(\Box\): the returned type is
a source image, and a normal source image is not that atom. Confluence
between two normal expressions therefore forces the sort and returned
type to be the literal \(\ast\). A logical annotation instead
gives its validity certificate. The annotation is already fully
normal by the input's annotation-normal hypothesis, so
\(\mathrm{Sort}\) accepts it and constructs the checked extension.
Apply the logical body induction there. Its target type is beta-equal
to \(\Box\), forcing the generated body sort to be \(\Box\).
The algorithm accepts the product's validity. This also proves the
asserted logical type classification, rather than assuming it.

For an object variable, generation makes its given type beta-equal
to the checked lookup; lookup is exactly the algorithm's normal
answer. For an object product, generate its annotation/body sorts.
Sort acceptance of its normal annotation is the case just proved.
The smaller object body induction supplies its normal type. As that
answer is a source image, it cannot be beta-equal to the atom
\(\Box\); equality to a generated sort forces the literal answer
\(\ast\). The product branch therefore accepts and returns its
actual formation sort \(\ast\), beta-equal to the original type.

For a lambda, generation gives the annotation sort, a body typing
\(TM:B\), and a well-formed product \(\Pi x:TA.B\). The smaller
annotation case checks its normal annotation. In this checked
extension the body induction returns a normal source type \(P\)
with \(TP=_\beta B\). The algorithm returns \([x:A]P\), which
is normal because both domain and body type are normal. Congruence
gives its target beta-equality to the generated product type and
therefore to the original type. This step does not assume that the
generated arbitrary target type \(B\) already has a source inverse.
The source preimage \(P\) was actually constructed by body inference.

For an application, generation gives
\(TM:\Pi x:A.B\), \(TN:A\), and result beta-equal to
\(B[TN/x]\). The two smaller inductions return normal source types
\(P,Q\), respectively beta-equal to this product and to \(A\).
The generated product is well formed by target type correctness.
Subject reduction and normalization preserve its product head, giving
\(\Pi x:\mathcal N(A).\mathcal N(B)\). Confluence with the
already normal \(TP\) forces \(P\) to have that literal product
head. Its source domain is normal. Confluence now forces the already
normal \(TQ\) to equal that domain literally. In indexed syntax
this is precisely the algorithm's argument comparison, not merely
an appeal to decidable conversion. The two successful premise
derivations sort the substituted result type, as in soundness.
Normalize it. Substitution congruence from the generated product's
components shows that the returned type is beta-equal to the
original application's result type. All application success and
failure branches are thus covered by generation. No source typed
equality reflection was used.

These inductions use only proper original syntax children. For
annotation checking, a smaller annotation's entire syntax is part
of the original input, not a generated type to be recursively
rechecked. Generated normal products are inspected and their domains
compared, not passed back to inference. The occasional full normalizer
calls have already typed CC inputs, so exact CC strong normalization
terminates them. All other checks, scope tests, raw redex tests,
lookups and literal comparisons are finite structural computations.
Consequently the three operations terminate on arbitrary finite
inputs, including unsuccessful branches, in a checked environment.

Now validate environments oldest-first. Each loop iteration invokes
these already total operations on the next annotation under its
checked older prefix. There are finitely many declarations. A success
builds its actual normal-variant validity derivation. Conversely take
any normal-variant valid environment. Its declaration annotations
are normal by transport, and its target context is valid. The target
annotation certificates, together with the foregoing completeness
under the already checked older prefix, force every next sort check
to accept. If the annotation is itself a normal form originally
introduced from another term, its target certificate comes from
subject reduction of that term, not an assumed source-preservation
lemma. Thus every valid environment is accepted. The same prefix argument
applies to any source-shaped context with normal annotations whose
encoding is CC-valid: its target annotation certificates suffice in
each checked older prefix by the stronger syntax completeness already
proved. This additional context-image assertion does not assume a
preexisting normal-variant derivation.

A normal-variant object typing now has an annotation-normal subject
by transport. Its CC typing certificate forces inference acceptance
by completeness. The returned normal type is beta-equal to its given
type, also normal by transport; confluence and indexed injectivity
make them literally identical. Logical validity similarly forces
acceptance. Combined with soundness and totality, rejection is
complete for the declared judgments, and supplied-type checking
is precisely comparison with the unique returned type.

For type correctness, use the target sort certificate of a derived
type \(P\) and its normality. If \(P\) is an object, apply the
stronger object completeness to \(P\) itself. Its normal inferred
type is equal to the target sort, which cannot be \(\Box\) by
the same normal-source-image argument, hence is \(\ast\).
Soundness supplies \(\Gamma\vdash_n P:\ast\). If \(P\) is
logical, the logical completeness supplies its source validity.
Finally the subject \(M\) has a target typing and hence a full
normal form. Reflection gives the source \(\mathcal N(M)\);
target subject reduction types it at the same \(TP\).
It is now fully normal and therefore annotation-normal. Stronger
completeness infers its source normal type, which confluence identifies
with \(P\). Soundness supplies
\(\Gamma\vdash_n\mathcal N(M):P\). This is a proved whole-normal-form
preservation result, not the false raw-step assertion.
\end{proof}
\end{annexitem}

\begin{annexitem}{annex:p7.ch88.check-normal-step}{thm:p07-ch88-check-normal-step}{The normal-annotation beta countercase}
\begin{proof}
In \(A:\ast,a:A\), the logical context \(\kappa=[x:A]\ast\)
is valid and normal. Abstraction from \(A:\ast\) in its
\(x:A\) extension gives \(K:\kappa\). Extend by
\(F:\kappa\); application gives \(Fa:\ast\), since its
result is the unchanged normal star. This type \(Fa\) is a normal
neutral application. Extend by \(y:Fa\), project \(y:Fa\), and
abstract to obtain \((\lambda y:Fa)y:[y:Fa]Fa\).
Abstract over \(F\), with normal domain \(\kappa\), obtaining
function type \([F:\kappa][y:Fa]Fa\), also normal. Apply to
\(K:\kappa\). The normal application result is
\(\mathcal N([y:Ka]Ka)=[y:A]A\), so the claimed typing of
\(V\) is an actual no-conversion derivation.

Its raw root contraction substitutes \(K\) for \(F\) and gives
\((\lambda y:Ka)y\). The domain \(Ka=((\lambda x:A)A)a\)
is a beta redex, so the result fails the necessary annotation-normal
invariant of every normal-variant typing. It is therefore not typable
there; this is an invariant argument, not merely checker rejection.
Reducing that domain yields \((\lambda y:A)y\), whose common
abstraction derivation has type \([y:A]A\). These examples preserve
all source object/logical classes and use no extra axiom or reduction.
They show why normal-form preservation and raw-step preservation must
be stated separately.
\end{proof}
\end{annexitem}
